\documentclass[11pt,a4paper]{article}
\usepackage[ngerman]{babel}
% Shared preamble of the English and German IT security guides (guide.tex, guide-de.tex).
\usepackage[utf8]{inputenc}
\usepackage[T1]{fontenc}
\usepackage{lmodern}
\usepackage{microtype}
\usepackage[margin=0.9in]{geometry}
\usepackage{booktabs}
\usepackage{tabularx}
\usepackage{array}
\usepackage{enumitem}
\usepackage{fontawesome5}
\usepackage{xcolor}
\usepackage{listings}
\usepackage{tcolorbox}
\usepackage[colorlinks=true,linkcolor=blue,citecolor=blue,urlcolor=blue]{hyperref}
\setlength{\emergencystretch}{2em}
\setlist{itemsep=2pt,topsep=3pt}

\definecolor{separblue}{HTML}{0550AE}
\definecolor{separbg}{HTML}{F6F8FA}
\definecolor{separborder}{HTML}{D0D7DE}
\definecolor{separwarn}{HTML}{FFF8C5}
\definecolor{separwarnborder}{HTML}{D4A72C}

\lstset{
  backgroundcolor=\color{separbg}, frame=single, rulecolor=\color{separborder},
  basicstyle=\ttfamily\small, breaklines=true, columns=fullflexible, keepspaces=true,
  showstringspaces=false, aboveskip=6pt, belowskip=6pt, xleftmargin=4pt, xrightmargin=4pt,
  upquote=true
}
\newtcolorbox{callout}[1][]{colback=separbg, colframe=separblue, boxrule=0.8pt, arc=3pt,
  left=8pt, right=8pt, top=6pt, bottom=6pt, #1}
\newtcolorbox{warnbox}[1][]{colback=separwarn, colframe=separwarnborder, boxrule=0.8pt, arc=3pt,
  left=8pt, right=8pt, top=6pt, bottom=6pt, #1}

% a fact the repository does not hold yet (listed in the build report)
\newcommand{\todo}[1]{\textcolor{red}{\textbf{[TODO:} #1\textbf{]}}}
\newcommand{\cmd}[1]{\texttt{#1}}
\newcommand{\contact}{\href{mailto:parsa@parsabe.com}{\texttt{parsa@parsabe.com}}}
\newcommand{\repo}{\href{https://github.com/parsabe/Separ}{\texttt{github.com/parsabe/Separ}}}
\newcommand{\releases}{\href{https://github.com/parsabe/Separ/releases}{\texttt{github.com/parsabe/Separ/releases}}}


\title{\Large\bfseries Separ v2.4 auf Linux-Rechnern der TU Bergakademie Freiberg\\[4pt]
\large Leitfaden für den IT-Sicherheitsbeauftragten, die Administratorinnen und Administratoren des
URZ und den Datenschutzbeauftragten}
\author{Parsa Besharat (studentisches Forschungsprojekt)\\ \contact \quad|\quad \repo}
\date{November 2026 \quad|\quad gilt für Separ 2.4.0}

\begin{document}
\maketitle

\begin{callout}
\textbf{Zum Gebrauch dieses Leitfadens.} Abschnitt~\ref{sec:summary} ist eine einseitige
Zusammenfassung für Entscheidungsträger. Die Abschnitte~\ref{sec:does}--\ref{sec:approvals}
erklären, was Separ tut und welche Freigaben nötig sind. Die
Abschnitte~\ref{sec:req}--\ref{sec:pilot} richten sich an die Administratorinnen und
Administratoren, die den Pilotbetrieb durchführen. Jeder Befehl und jede Option in diesem
Leitfaden stammt aus \texttt{docs/it-facts.md} im Repository, das die genauen Hilfetexte von
Separ 2.4.0 festhält. Netzbereiche, Namen und Adressen in Beispielen sind Platzhalter wie
\texttt{<your campus range>}. Meldungen und Befehlsausgaben von Separ sind englisch und werden
hier im Original wiedergegeben.
\end{callout}

\tableofcontents
\newpage

% ============================================================================
\section{Zusammenfassung für Entscheidungsträger}
\label{sec:summary}

\paragraph{Was Separ ist.} Ein studentisches Forschungsprojekt (Version 2.4.0), das Vorgänge auf
einem Linux-Rechner betrachtet und drei Fragen stellt: Kam die Eingabe von außen (eine Webseite,
eine E-Mail, ein Download, die Eingabe eines KI-Agenten, eine Anfrage an einen Webdienst)? Tut das
Programm etwas, das seine Rolle nie tut (ein Webserver startet eine Shell)? Versucht es etwas
Gefährliches (Code herunterladen und ausführen, sich verstecken, sich beim Systemstart
eintragen, Passwörter lesen)? Für Befehle, die über Separ geleitet werden, entscheidet Separ
\emph{bevor} sie laufen: zulassen, anhalten (nachfragen) oder blockieren. Für Sensordaten, die
eine Administratorin oder ein Administrator anbindet, entscheidet es im Nachhinein. Separ enthält
kein maschinelles Lernen und sendet nichts ins Internet.

\paragraph{Worum gebeten wird.} Ein \textbf{reiner Beobachtungs-Pilot auf 3 bis 5 Linux-Rechnern
für etwa zwei Wochen}. Im Beobachtungsmodus (\cmd{observe}) blockiert Separ nichts: Es hält nur
fest, was es anhalten oder blockieren \emph{würde}, damit Fehlalarme bei echter Arbeit gezählt
werden können, bevor jemand eine Durchsetzung erwägt.

\paragraph{Was es aufzeichnet.} Auf jedem Rechner eine Zeile je beobachtetem Ereignis und je
Entscheidung. Darin sind Befehle, Webadressen, Dateipfade und Dateiinhalte durch verschlüsselte
Fingerabdrücke ersetzt (HMAC-SHA256 mit einem zufälligen Schlüssel je Rechner). Einige Kennungen
speichert es jedoch im Klartext: Anmeldenamen, von Sensoren gemeldete Rechnernamen und die Namen
in Quarantäne verschobener Dateien (Abschnitt~\ref{sec:facts}). Nichts verlässt den Rechner, es
sei denn, eine Administratorin oder ein Administrator richtet einen Benachrichtigungskanal ein;
diese Nachrichten nennen den Rechner nur unter einem Pseudonym und enthalten keinen Benutzernamen.

\paragraph{Was es nicht kann.}
\begin{itemize}
  \item Es sieht nur, was ihm zugeleitet wird: Befehle über \cmd{separ guard}, die
        Browsererweiterung, einen überwachten Ordner und angebundene Sensoren. Ein Befehl, der ohne
        \cmd{separ guard} in ein Terminal eingefügt wird, oder ein Skript, das über seinen Namen
        gestartet wird, wird vor der Ausführung nicht gesehen.
  \item Im Blockiermodus blockiert es auch legitime Installationsanweisungen, die von einer
        Webseite eingefügt werden (im Test wurden 5 von 10 alltäglichen eingefügten Befehlen
        blockiert). Deshalb ist der Pilot ein reiner Beobachtungsbetrieb.
  \item Echte Prozesse werden nur eingedämmt, wenn der Systemdienst als root mit \cmd{--enforce}
        läuft; sonst wird die Eindämmung nur simuliert. root, und bei Einrichtungen je Benutzer
        derselbe Benutzer, kann Separ beenden.
  \item Es wurde nicht an echten Benutzern oder Servern gemessen. CPU- und Speicherbedarf sind
        unbekannt und sollen im Pilotbetrieb gemessen werden.
\end{itemize}

\paragraph{Benötigte Entscheidungen.} Die Freigabe durch den IT-Sicherheitsbeauftragten, eine
datenschutzrechtliche Entscheidung des Datenschutzbeauftragten und, bei Rechnern von
Beschäftigten oder Servern, an denen sich Beschäftigte anmelden, die Beteiligung des
Personalrats (Abschnitt~\ref{sec:approvals}).

% ============================================================================
\section{Was Separ tut und was nicht}
\label{sec:does}

\subsection{Die drei Fragen}
\begin{description}
  \item[Herkunft (Origin).] Kam die auslösende Eingabe von außerhalb der Vertrauensgrenze? Inhalte
    von einer Webseite, aus einer E-Mail, aus einem Download, von einem Wechseldatenträger, eine
    Anfrage aus dem Internet oder die Eingabe eines KI-Agenten erhalten eine \emph{Markierung}
    (Label), die über 30 Minuten verblasst, sofern der Vorgang nicht gefährlich wird.
  \item[Vertrag (Contract).] Handelt das Programm außerhalb seiner Rolle? Jede Rolle (interaktive
    Shell, Webserver-Prozess, Paketinstallation, \ldots) hat einen signierten Vertrag darüber, was
    sie darf.
  \item[Wirkung (Effect).] Versucht der Vorgang etwas Gefährliches: von einem unbekannten Ort
    nachladen, sich verstecken, sich dauerhaft einnisten, Rechte ausweiten, sich im Netz seitwärts
    bewegen, Zugangsdaten lesen, Sicherheitswerkzeuge manipulieren, viele Dateien auf einmal
    ändern?
\end{description}
Separ addiert die Belege der drei Fragen zu einem Wert. Über einer Schwelle hält es an (fragt
nach), über einer höheren blockiert es; ohne Vertragsverletzung blockiert es nie allein auf
Verdacht. Einige Ereignisse blockieren immer: das Lesen eines Köder-Passworts (Canary), der
Zugriff eines markierten Programms auf ein geschütztes Gut und Versuche, Separ zu beenden oder zu
verändern.

\subsection{Wo es hinsieht}
\begin{tabularx}{\textwidth}{@{}lX@{}}
\toprule
\textbf{Kontrollpunkt} & \textbf{Wann es entscheidet} \\
\midrule
\cmd{separ guard -- <command>} & bevor der Befehl läuft \\
Browsererweiterung & wenn eine Seite in die Zwischenablage schreibt, ein Download beginnt,
  optional beim Öffnen einer Seite \\
Überwachte Sandbox und Web-Root & wenn eine Datei erscheint \\
Angebundene Sensoren (Tetragon, Falco, sshd, Shibboleth, \ldots) & nach dem Ereignis, aus deren
  Protokollen \\
Rootkit-Prüfung (\cmd{separ scan-rootkit}) & regelmäßig, nur lesend \\
\bottomrule
\end{tabularx}

\subsection{Ehrliche Grenzen}
\begin{itemize}
  \item \textbf{Eingefügte Installationsbefehle werden im Blockiermodus blockiert.} Ein von einer
    Webseite eingefügter Befehl, der sichtbar herunterlädt, Rechte ausweitet oder in eine Shell
    weiterleitet, scheitert an allen drei Fragen. Im Testsatz wurden so \cmd{git clone}, das
    \cmd{rustup}-Installationsprogramm, \cmd{sudo apt update} und \cmd{wget} blockiert. Vertrauenswürdige
    Ziele einzutragen ändert daran nichts.
  \item \textbf{Getippte Befehle, die nachladen oder Rechte ausweiten, werden angehalten}
    (nachgefragt), nie blockiert.
  \item \textbf{Nicht gesehen:} ein Einfügen ins Terminal ohne \cmd{separ guard} oder die
    Erweiterung, ein über seinen Namen gestartetes Skript, ein kompiliertes Programm, alles nicht
    Angebundene.
  \item \textbf{Beenden von Separ:} root kann den Systemdienst beenden; ein Benutzer kann seinen
    eigenen Dienst je Benutzer beenden. Separ meldet sein Beenden, und eine externe
    Zustandsprüfung bemerkt einen fehlenden Heartbeat, aber beides verhindert es nicht.
  \item \textbf{Eindämmung} wird nur simuliert, es sei denn, der Systemdienst läuft als root mit
    \cmd{--enforce} (Abschnitt~\ref{sec:extras}); auch dann werden nur Prozesse angefasst, die ein
    Sensor mit Prozess-ID gemeldet hat.
\end{itemize}

\subsection{Separ und Falco}
Falco beobachtet Systemaufrufe und meldet \emph{nach} einer Aktion; Separ beurteilt die ihm
zugeleiteten Befehle, \emph{bevor} sie laufen. Beide erfüllen unterschiedliche Aufgaben und können
zusammen betrieben werden: Separ kann die Meldungen von Falco als Sensor einlesen
(Abschnitt~\ref{sec:extras}). In einem Labortest mit 22 harmlosen Angriffs-Nachbildungen meldete
Falco 0.45.2 mit seinen 25 Standardregeln nichts; mit 81 Regeln (einschließlich seiner
experimentellen Regelsätze) meldete es 3 davon und 1 von 6 harmlosen Befehlen. Separ hielt 21 an
oder blockierte sie und keinen der harmlosen Befehle. \textbf{Dieser Test begünstigt Separ}: Die
Nachbildungen stammen vom Autor von Separ, die Standardregeln von Falco zielen vor allem auf
Container, die Nachbildungen erreichen kein echtes Netz, und alltägliche Befehle waren nicht Teil
des Tests. Er zeigt nicht, wie sich eines der beiden Werkzeuge gegen echte Angriffe schlägt
(Einzelheiten: \texttt{docs/v2.4-evaluation.md} im Repository).

% ============================================================================
\section{Freigaben vor der Installation}
\label{sec:approvals}

\begin{enumerate}
  \item \textbf{IT-Sicherheitsbeauftragter}: gibt den Pilotbetrieb, die Rechner und die Personen
    mit Zugriff auf die Protokolle frei.
  \item \textbf{Datenschutzbeauftragter}: entscheidet, ob eine Datenschutz-Folgenabschätzung nach
    Art.~35 DSGVO erforderlich ist, und nimmt die Verarbeitung in das Verzeichnis von
    Verarbeitungstätigkeiten auf (Art.~30 DSGVO). Das Datenblatt in Abschnitt~\ref{sec:facts}
    listet auf, was gespeichert wird.
  \item \textbf{Personalrat}: bei Rechnern von Beschäftigten und bei Servern, an denen sich
    Beschäftigte anmelden, weil Separ dort Benutzeraktivität aufzeichnet.
  \item \textbf{Schriftliche Information der Benutzerinnen und Benutzer} der Pilotrechner vor
    Beginn: was aufgezeichnet wird, wie lange, wer es sehen kann und an wen man sich wendet.
  \item \textbf{Private Geräte}: nur mit ausdrücklicher Einwilligung der Eigentümerin oder des
    Eigentümers.
\end{enumerate}

\subsection{Datenblatt Datenschutz}
\label{sec:facts}
Aus einem Laborlauf von Separ 2.4.0 (\texttt{docs/it-facts.md}, Abschnitt 4).

\noindent{\small
\begin{tabularx}{\textwidth}{@{}>{\raggedright\arraybackslash}p{0.3\textwidth}X>{\raggedright\arraybackslash}p{0.22\textwidth}@{}}
\toprule
\textbf{Datei} & \textbf{Daten} & \textbf{Form} \\
\midrule
\cmd{logs/events.jsonl} & Befehl, URL, Pfad, Dateiinhalt, Anmeldezeile & verschlüsselter Fingerabdruck \\
 & Rechner bei überwachten Befehlen, Browser, Sandbox & Pseudonym \\
 & Rechner bei Ereignissen angebundener Sensoren & Rechnername im \textbf{Klartext} \\
 & Anmeldename (\cmd{actor.user}); Konto bei Anmelde-Sensoren & \textbf{Klartext} \\
 & Quelladresse einer Anmeldung & Klasse oder verschlüsselter Fingerabdruck \\
\cmd{logs/verdicts.jsonl} & Entscheidung, Wert, Erklärung; Benutzernamen & \textbf{Benutzernamen im Klartext} \\
\cmd{logs/notifications.jsonl}, Webhook & Rechner, feste Sätze, Entscheidung & Pseudonym, kein Benutzer \\
\cmd{logs/barrier.log}, Dashboard & Namen von Dateien in Quarantäne; Pfad des Arbeitsverzeichnisses & \textbf{Klartext} \\
\cmd{heartbeat.json} & Rechner, Prozess-ID des Dienstes & Pseudonym \\
Köder-Vorschläge (nur mit Ködern) & Adresse einer externen / internen Quelle & Klartext / versiegelt \\
\bottomrule
\end{tabularx}}

\begin{warnbox}
\textbf{Rechte der Protokolldateien.} Der Systemdienst legt seine Protokolle mit der
Standard-umask in \cmd{/var/lib/separ} (Modus 0755) an, daher können in Separ 2.4.0
\textbf{alle lokalen Benutzer sie lesen}. Schränken Sie dies vor dem Pilotbetrieb mit einer
Drop-in-Datei ein: \cmd{sudo systemctl edit separ}, die beiden Zeilen \cmd{[Service]} und
\cmd{UMask=0077} hinzufügen, dann \cmd{sudo systemctl restart separ} (bereits angelegte Dateien
behalten ihre Rechte). Separ rotiert seine Protokolle nicht; löschen oder archivieren Sie sie am
Ende des Pilotbetriebs wie mit dem Datenschutzbeauftragten vereinbart.
\end{warnbox}

Aufbewahrung, Zugriff und Löschung legt der Datenschutzbeauftragte fest; Separ selbst löscht
nichts. Das Dashboard zeigt die letzten Ereignisse, auf einem Server mit angebundenen Sensoren also
auch Kontonamen anderer Benutzer; geben Sie den Zugang zum Dashboard nur den in der Freigabe
genannten Personen.

% ============================================================================
\section{Mindestanforderungen}
\label{sec:req}
\begin{itemize}
  \item Linux auf x86\_64 (\cmd{uname -m} zeigt \cmd{x86\_64}) mit systemd; Ubuntu 24.04,
    Debian 13 oder ein System der RHEL-Familie 9 mit Python 3.12 (wird mit dem Paket installiert,
    siehe Abschnitt~\ref{sec:install}). Debian 12 wird nicht unterstützt (Python 3.11).
  \item Einmalig root-Rechte, um das Paket und die signierte Richtlinie zu installieren.
  \item Kein Internetzugang bei der Installation nötig (das Paket enthält seine
    Python"=Abhängigkeiten), keine eingehenden Verbindungen: der Dienst lauscht nur
    auf \cmd{127.0.0.1} (standardmäßig auf Port 8765).
  \item Kein Kernelmodul, kein Cloud-Konto, keine Lizenz.
  \item Hardware: \textbf{nicht gemessen}. CPU-, Speicher- und Plattenbedarf sollen im
    Pilotbetrieb gemessen werden. Eine vollständige Rootkit-Prüfung dauerte auf einem
    Entwicklungs-Laptop etwa 26 Sekunden.
\end{itemize}

% ============================================================================
\section{Rollen und Schlüsselverwahrung}
\label{sec:roles}
\begin{description}
  \item[Richtlinienverantwortliche] (eine Administratorin oder ein Administrator des URZ und eine
    Vertretung). Verwahrt den privaten Signaturschlüssel \cmd{admin.key} auf einem
    \textbf{Administrationsrechner} (Standard \cmd{\textasciitilde/.separ-admin/admin.key}, Modus
    0600), \textbf{nie auf einem geschützten Rechner}. Hält eine Offline-Sicherung vor (zum
    Beispiel einen verschlüsselten USB-Stick in einem Tresor). Signiert jede Änderung der
    Richtlinie.
  \item[Bearbeitende (Responder).] Lesen Benachrichtigungen und das Dashboard, beantworten
    Anhaltungen auf Servern und erheben die Kennzahlen des Pilotbetriebs.
\end{description}
Die geschützten Rechner erhalten nur \cmd{policy.yaml} und \cmd{admin.pub}; der Dienst prüft die
Signatur und kann selbst nichts signieren.

\paragraph{Wenn der Schlüssel verloren geht.} Bereits installierte Richtlinien funktionieren
weiter, aber keine Änderung kann signiert werden. Erzeugen Sie auf dem Administrationsrechner einen
neuen Schlüssel und eine neue Richtlinie (Abschnitt~\ref{sec:policy}, mit einem neuen Pfad für
\cmd{--key} und einem neuen \cmd{--config-dir}), installieren Sie die neue \cmd{admin.pub} und
\cmd{policy.yaml} auf jedem Rechner und laden Sie neu.

\paragraph{Wenn der Schlüssel gestohlen wird.} Wer ihn besitzt, kann eine Richtlinie signieren,
die Separ in den Beobachtungsmodus schaltet oder alles erlaubt. Separ hat keine Sperrliste:
Ersetzen Sie \cmd{admin.pub} und \cmd{policy.yaml} auf allen Rechnern gleichzeitig (wie oben),
vergleichen Sie den Digest, den \cmd{separ policy show} auf jedem Rechner ausgibt, mit dem von
Ihnen signierten, und melden Sie den Vorfall dem IT-Sicherheitsbeauftragten.

% ============================================================================
\section{Installation}
\label{sec:install}
Laden Sie das Paket und \cmd{SHA256SUMS} der Version von \releases{} herunter, prüfen und
installieren Sie es:
\begin{lstlisting}
BASE=https://github.com/parsabe/Separ/releases/download/v2.4
curl -fLO $BASE/SHA256SUMS
# Ubuntu 24.04 / Debian 13
curl -fLO $BASE/separ_2.4.0_amd64.deb
sha256sum --ignore-missing -c SHA256SUMS
sudo apt install ./separ_2.4.0_amd64.deb
# RHEL-family 9
curl -fLO $BASE/separ-2.4.0-1.x86_64.rpm
sha256sum --ignore-missing -c SHA256SUMS
sudo dnf install python3.12 ./separ-2.4.0-1.x86_64.rpm
\end{lstlisting}
Die Version wird von der CI des Projekts auf Ubuntu 24.04, Debian 13 und Rocky Linux 9 gebaut und
installationsgetestet \todo{bestätigen, dass die Version v2.4 mit ihren Paketen veröffentlicht
ist}.

\paragraph{Wohin die Dateien kommen.} \cmd{/usr/bin/separ}; eine eigene Python-Umgebung in
\path{/opt/separ/venv}, gebaut aus den Wheels in \path{/opt/separ/wheels}; Vorlagen in
\cmd{/etc/separ/} (\cmd{ingest.yaml}, Beispiele, \cmd{mappings/}); die System-Unit
\path{/usr/lib/systemd/system/separ.service} und die Benutzer-Unit
\path{/usr/lib/systemd/user/separ-user.service}, beide nicht aktiviert; die Browsererweiterung und
ihre Richtlinienvorlagen in \cmd{/usr/share/separ/}. \textbf{Es werden weder Richtlinie noch
Schlüssel installiert.} Zustand und Protokolle liegen in \cmd{/var/lib/separ}.

\paragraph{Prüfung.}
\begin{lstlisting}
separ --version            # separ 2.4.0
separ health --json        # "not_running" (exit code 1) until the service is started
\end{lstlisting}

% ============================================================================
\section{Arbeitsablauf für die Richtlinie}
\label{sec:policy}

\subsection{Richtlinie auf dem Administrationsrechner erstellen}
\begin{lstlisting}
separ setup --server --mode observe --config-dir ./tubaf-policy \
  --campus-cidr <your campus range> --vpn-cidr <your VPN range> \
  --wlan-cidr <your WLAN range> --mail-domain <your mail domain> \
  --key ~/.separ-admin/admin.key
\end{lstlisting}
Verwenden Sie für Arbeitsplatzrechner \cmd{--laptop} statt \cmd{--server}. Mit \cmd{--server} ist
\cmd{--campus-cidr} Pflicht; Separ hat keinen Campus-Bereich fest eingebaut. Der Befehl gibt genau
aus, was er eingerichtet hat, und schreibt \cmd{policy.yaml} und \cmd{admin.pub} nach
\cmd{./tubaf-policy}; der Schlüssel wird angelegt, falls er noch nicht existiert.

\subsection{Auf einem geschützten Rechner installieren}
\begin{lstlisting}
scp tubaf-policy/policy.yaml tubaf-policy/admin.pub <host>:
sudo install -o root -g root -m 0644 policy.yaml /etc/separ/policy.yaml
sudo install -o root -g root -m 0644 admin.pub /etc/separ/admin.pub
separ policy show --policy /etc/separ/policy.yaml     # signature, mode, digest
sudo systemctl enable --now separ.service
\end{lstlisting}
Der Dienst startet nicht ohne signierte Richtlinie und weist eine Richtlinie oder einen Schlüssel
zurück, die ein symbolischer Link sind, nicht root gehören oder für Gruppe oder andere schreibbar
sind. Dienste je Benutzer auf vom URZ verwalteten Rechnern verwenden \cmd{/etc/separ}, sobald dort
eine Richtlinie liegt, nach denselben Regeln.

\subsection{Modus oder Richtlinie ändern}
Auf dem Administrationsrechner, dann wie oben kopieren und installieren, dann neu laden:
\begin{lstlisting}
separ policy set-mode warn --policy tubaf-policy/policy.yaml --key ~/.separ-admin/admin.key
separ policy sign --policy tubaf-policy/policy.yaml --key ~/.separ-admin/admin.key  # after editing
sudo separ reload --workdir /var/lib/separ        # or: sudo systemctl reload separ
\end{lstlisting}
Das Neuladen liest die Datei nur neu ein und prüft sie; eine zurückgewiesene Datei lässt die
letzte gültige Richtlinie in Kraft, und \cmd{separ health} meldet \emph{Policy rejected}.

% ============================================================================
\section{Plan für den Pilotbetrieb}
\label{sec:pilot}

\paragraph{Aufbau.} 3--5 Linux-Rechner, Beobachtungsmodus, etwa zwei Wochen, Freigaben liegen vor,
Benutzerinnen und Benutzer sind informiert, Rechte der Protokolle eingeschränkt
(Abschnitt~\ref{sec:facts}). Wer eingefügte Befehle beurteilen lassen möchte, startet sie über
\cmd{separ guard -- <command>}; mit der (optionalen) Browsererweiterung werden Schreibvorgänge von
Seiten in die Zwischenablage auch ohne dies beurteilt.

\paragraph{Wöchentlich, je Rechner.}
\begin{lstlisting}
sudo grep -c 'Would hold (not enforced)'  /var/lib/separ/logs/notifications.jsonl
sudo grep -c 'Would block (not enforced)' /var/lib/separ/logs/notifications.jsonl
\end{lstlisting}
Halten Sie fest: die Zahl der Vorgänge, die angehalten oder blockiert worden wären; für jeden, was
es war (Benachrichtigung und Dashboard erklären es in Sätzen; Separ speichert den Befehlstext nicht,
fragen Sie daher die Benutzerin oder den Benutzer, was zu dieser Zeit ausgeführt wurde); und ob es
ein Fehlalarm war. Notieren Sie außerdem CPU- und Speicherbedarf des Prozesses \cmd{separ} und die
Größe von \cmd{/var/lib/separ/logs}.

\paragraph{Das Dashboard.} \cmd{sudo separ gui --workdir /var/lib/separ --no-browser} gibt die
lokale Adresse mit ihrem Token aus; öffnen Sie sie auf dem Rechner selbst (oder über einen
SSH-Tunnel zu Port 8765). Es zeigt jeden Vorgang mit seinen drei Fragen, seinem Wert und seiner
Erklärung.

\paragraph{Kriterien für den nächsten Schritt} (vor dem Pilotbetrieb mit dem
IT-Sicherheitsbeauftragten festzulegen):
\begin{itemize}
  \item für den \textbf{Warnmodus} (\cmd{warn}; Benutzer sehen einen Hinweis, nichts wird
    gestoppt): keine Blockierung bei Routinearbeit, die sich nicht über die Richtlinie beheben
    lässt, und höchstens etwa eine Anhaltung je Benutzer und Woche;
  \item für den \textbf{Blockiermodus} (\cmd{block}): ein weiterer Zeitraum im Warnmodus mit
    denselben Zielen, eine Entscheidung über eingefügte Installationsbefehle (sie werden
    blockiert) und ein vereinbarter Ablauf für die Reaktion.
\end{itemize}

% ============================================================================
\section{Optionale Erweiterungen}
\label{sec:extras}
Jede Erweiterung hat eigene Voraussetzungen und wird für den Pilotbetrieb nicht benötigt.
\begin{description}
  \item[Browsererweiterung.] Erfordert Firefox ESR 128+ oder Chromium/Chrome 121+ und einen Dienst
    je Benutzer. Als root (für Chromium zusätzlich \cmd{--chromium-extension-id <id>}):
\begin{lstlisting}
sudo separ browser-install --system
\end{lstlisting}
    anschließend die Erweiterung mit den Vorlagen für Unternehmensrichtlinien in
    \path{/usr/share/separ/browser/} verteilen.
  \item[Echte Durchsetzung.] Erfordert den Systemdienst als root, cgroup v2 und \cmd{nft} für die
    Netztrennung sowie einen Sensor, der Prozess-IDs meldet (Tetragon oder Falco). Ergänzen Sie
    \cmd{--enforce} in der Zeile \cmd{ExecStart} der Unit mit \cmd{sudo systemctl edit separ}.
    \cmd{separ health} listet danach jede Eindämmungsmaßnahme als verfügbar oder nicht.
  \item[Kernel-Sensoren] (Tetragon, Falco). Installieren Sie den Sensor, aktivieren Sie seine
    Zuordnung (Mapping) in \path{/etc/separ/ingest.yaml} und testen Sie zuerst eine Stichprobe:
\begin{lstlisting}
separ ingest list
separ ingest check /etc/separ/mappings/<mapping>.yaml --input <sample> --judge
\end{lstlisting}
  \item[Anmelde- und Mail-Sensoren] (sshd, Shibboleth IdP, Keycloak, Postfix/Dovecot,
    M365-Audit). Ihre Regeln beginnen im Beobachtungsmodus, und ein Vorgang, den nur sie erzeugt
    haben, wird höchstens angehalten. Sie zeichnen Kontonamen im Klartext auf; beziehen Sie sie in
    die datenschutzrechtliche Entscheidung ein.
  \item[Benachrichtigungen.] Optionen von \cmd{separ serve} (in \cmd{ExecStart} ergänzen);
    \cmd{--notify-on} wählt die Kategorien. Nachrichten nennen den Rechner nur unter einem
    Pseudonym.
\begin{lstlisting}
--notify-webhook <url>          # a Mattermost or Slack incoming webhook
--notify-command "mail -s Separ admin@example.org"
\end{lstlisting}
\end{description}

% ============================================================================
\section{Betrieb}
\label{sec:ops}

\paragraph{Zustandsprüfung und Überwachung.} \cmd{separ health} liest die Heartbeat-Datei:
Rückgabewert 0 gesund, 1 nicht laufend oder veraltet, 2 laufend mit einem Problem. Ein
cron-Eintrag (\cmd{crontab -e} als root) benachrichtigt einmal je Ausfall und einmal bei der
Wiederherstellung:
\begin{lstlisting}
*/5 * * * * separ health --workdir /var/lib/separ --notify-command "mail -s Separ admin@example.org"
\end{lstlisting}

\paragraph{Rootkit-Prüfungen.} Der Dienst prüft stündlich (\cmd{--rootkit-interval}, Standard
3600~s) und durchsucht dabei nur die bisher vergebenen Prozess-IDs. Eine vollständige manuelle
Prüfung: \cmd{sudo separ scan-rootkit} (etwa 26~s auf einem Entwicklungs-Laptop;
\cmd{--pid-sweep bounded} ist schneller). Messen Sie die Rate vorübergehender Abweichungen einmal
auf einem sauberen Rechner:
\begin{lstlisting}
sudo separ scan-rootkit --calibrate 20 --workdir /var/lib/separ
\end{lstlisting}
Rückgabewerte: 0 nichts zu prüfen, 1 Befunde zu prüfen, 2 Befunde auf Blockierniveau.

\paragraph{Aktualisieren.} Installieren Sie das Paket der neuen Version wie in
Abschnitt~\ref{sec:install}, dann \cmd{sudo systemctl restart separ}. Die Richtlinie in
\cmd{/etc/separ} und der Zustand in \cmd{/var/lib/separ} bleiben erhalten.

\paragraph{Deinstallieren und Zurückrollen.} \cmd{sudo apt remove separ} oder \cmd{sudo dnf remove
separ} entfernt das Programm, die eigene Python-Umgebung und die vom Paket in \cmd{/etc/separ}
installierten Dateien; Richtlinie, öffentlicher Schlüssel und \cmd{/var/lib/separ} (Protokolle)
bleiben und müssen wie vereinbart gelöscht oder archiviert werden. Zum Zurückrollen entfernen Sie
das Paket und installieren die Paketdatei der vorherigen Version.

\paragraph{Fehlerbehebung.}
\begin{description}
  \item[\glqq no signed policy at \ldots: run \cmd{separ setup} first\grqq]
    Legen Sie die Dateien \cmd{policy.yaml} und \cmd{admin.pub} in \cmd{/etc/separ} ab (Abschnitt~\ref{sec:policy}).
  \item[\glqq refusing \ldots: writable by group or others\grqq{} / \glqq not root\grqq] Installieren
    Sie die Dateien erneut mit \cmd{install -o root -g root -m 0644}; auch der Ordner
    \cmd{/etc/separ} muss root gehören und darf für Gruppe und andere nicht schreibbar sein.
  \item[Policy rejected (\cmd{unsigned}, \cmd{bad\_signature})] Die Datei wurde ohne Signatur
    bearbeitet oder mit einem anderen Schlüssel signiert. Signieren Sie sie auf dem
    Administrationsrechner und installieren Sie sie erneut.
  \item[Degraded protection] Listet auf, was Separ auf diesem Rechner nicht sehen kann (zum
    Beispiel kein Kernel-Sensor, simulierte Eindämmung). Im Pilotbetrieb normal; eine Information,
    kein Fehler.
\end{description}

% ============================================================================
\section{Häufige Fragen von Benutzern und Service-Desk}
\label{sec:faq}
\begin{description}
  \item[Warum wurde mein Befehl angehalten?] Er kam von außen (zum Beispiel von einer Webseite
    eingefügt) oder tat etwas, das sein Programm normalerweise nicht tut, und sah danach aus, als
    könne er nachladen, sich verstecken, sich einnisten oder an Zugangsdaten gelangen. Im
    Beobachtungsmodus wurde nichts gestoppt; im Blockiermodus lief ein angehaltener Befehl, der über
    \cmd{separ guard} gestartet wurde, nicht.
  \item[Warum wurde ein normaler Installationsbefehl blockiert?] Von einer Webseite eingefügte
    Befehle, die Code herunterladen und ausführen, sehen genauso aus wie die Köder
    \glqq reparieren Sie Ihren Browser\grqq. Von Hand eingetippt werden sie nur angehalten.
    Während des Pilotbetriebs wird nichts blockiert.
  \item[Wie melde ich einen Fehlalarm?] Nennen Sie dem Service-Desk die Uhrzeit, den Rechner und
    was Sie gerade getan haben; die Bearbeitenden finden den Vorgang im Dashboard. Die
    Richtlinienverantwortlichen können eine eng gefasste, befristete Ausnahme in die Richtlinie
    aufnehmen und sie signieren.
  \item[Liest Separ meine Befehle oder Dateien?] Es speichert keinen Befehlstext, keine Adressen und
    keine Dateiinhalte, nur verschlüsselte Fingerabdrücke. Es zeichnet jedoch Ihren Anmeldenamen und
    die Uhrzeit jedes Ereignisses auf.
  \item[An wen wende ich mich?] An den Service-Desk des URZ; für das Projekt: \contact.
\end{description}

\end{document}
